\chapter{NMDA-Receptor Antibody Test}

\section*{What Is This Test?}
NMDA-receptor antibody testing detects immunoglobulin G antibodies against the NR1 subunit of the N-methyl-D-aspartate receptor. These antibodies are associated with anti-NMDA-receptor encephalitis, an autoimmune disorder that can cause psychiatric symptoms, seizures, movement abnormalities, impaired consciousness, and autonomic instability.

\section*{Alternative Names}
Anti-NMDAR Antibody; NMDA-Receptor Encephalitis Antibody; NMDAR IgG; GluN1 Antibody

\section*{How It Is Performed}
Antibodies are tested in cerebrospinal fluid and serum, commonly by cell-based immunofluorescence or another validated antibody assay. Cerebrospinal-fluid testing is particularly important because it is more specific for central nervous system disease; paired serum and cerebrospinal-fluid testing may improve interpretation.

\section*{Why It Is Ordered}
\begin{itemize}
  \item Evaluation of suspected autoimmune encephalitis with new psychiatric, cognitive, seizure, movement, or autonomic symptoms
  \item Investigation of unexplained encephalitis, especially in children, adolescents, and young adults
  \item Supporting diagnosis and follow-up of known anti-NMDA-receptor encephalitis in a specialist setting
\end{itemize}

\section*{Interpretation and Limitations}
A positive cerebrospinal-fluid result strongly supports anti-NMDA-receptor encephalitis in the appropriate clinical setting. Low-level serum-only reactivity can be nonspecific, and a negative result does not exclude autoimmune encephalitis caused by another antibody or a seronegative disorder. Results must be interpreted with neurologic examination, cerebrospinal-fluid studies, electroencephalography, and other investigations; brain imaging is separate from this test.
\section*{Regulatory Status and Authorization Date}
Regulatory status applies to the specific assay, instrument, manufacturer, and intended use rather than to the test name alone. No single approval date can be assigned to this test category. For a commercial assay, verify the current product in the relevant regulator's database: the U.S. Food and Drug Administration (FDA) clearance, approval, or authorization; India's Central Drugs Standard Control Organisation (CDSCO) licence or permission; the European Union In Vitro Diagnostic Regulation (EU IVDR) CE certificate issued through the applicable conformity-assessment route; or the United Kingdom Medicines and Healthcare products Regulatory Agency (MHRA) registration and UKCA/CE status. Laboratory-developed tests may not have an individual FDA, CDSCO, EU IVDR, or MHRA approval date.
\clearpage
